\section{Evaluation Methodology}
\label{sec:evaluation}

This report presents a compact table of selected public results
(Section~\ref{sec:results}, Table~\ref{tab:results}), each labeled with its
protocol and evidence tier. This section documents the \emph{methodology} behind
those labels --- how Argus separates kinds of evidence, and what measurement
discipline it applies --- so that every number can be placed in the correct tier.

\subsection{Evidence tiers}

Argus classifies every quantitative statement into one of five tiers. The tier,
not the magnitude, determines what the number is allowed to claim.

\begin{table}[h]
\centering
\small
\renewcommand{\arraystretch}{1.35}
\begin{tabularx}{\linewidth}{@{}p{3.2cm} X p{3.0cm}@{}}
\toprule
\textbf{Tier} & \textbf{What it means} & \textbf{What it may claim} \\
\midrule
Reproducible release evidence &
An in-repository artifact --- raw scored rows, logs, manifests, and a build script ---
that a third party can regenerate from a named commit. &
An Argus result for that specific run, on stated hardware. \\
Public website snapshot &
A published result card on the live public record~\cite{argus_results}, quoted
verbatim with its protocol, without an in-repository reproduction artifact yet. &
A scoped public claim under its stated protocol; linked to the evidence bundle. \\
Operator-observed preliminary &
A number seen during a live run but not yet packaged as a reproducible artifact
(e.g.\ a mid-run log). &
A preliminary observation only; not a result, not a claim. \\
External baseline &
A published third-party target from other work, used for comparison. &
A reference target; explicitly \emph{not} an Argus achievement. \\
Ongoing / no public claim &
A program actively running without certified, reproducible completion. &
Status ``ongoing''; no numeric Argus claim. \\
\bottomrule
\end{tabularx}
\caption{The five evidence tiers. A number's tier is fixed by its provenance, and
labels travel with the number wherever it appears.}
\label{tab:tiers}
\end{table}

\subsection{Measurement discipline}

Three rules govern how measurements are produced and reported.

\begin{description}[leftmargin=1.4em, style=nextline]
  \item[Same-hardware, same-harness baselines.] A comparison is only fair when the
    baseline is reproduced on the same hardware and measurement harness as the
    Argus run, with caveats stated explicitly; cross-hardware numbers are labeled,
    not silently compared.
  \item[Anti-cheat by construction.] Evaluation uses real public benchmark data.
    Where a metric could be gamed by memorizing a fixed evaluation input
    distribution, the input is randomized per evaluation so a solution cannot
    hardcode statistics of a known distribution --- a fast-because-faked kernel is
    designed to be impossible to pass off as genuinely fast. Every round's
    evidence, verdict, and skill update is persisted for audit.
  \item[Reviewer-certified completion.] A program is ``complete'' only when the
    Reviewer certifies it against the full-pipeline stage checklist. There is no
    hardcoded completion gate and no independent validator; the structured verdict
    is the single source of truth.
\end{description}

\subsection{Evidence map}

The claims in this report are grounded in the current \code{argus-skill}
repository~\cite{argus_repo} rather than in ephemeral run directories. Readers who
want to verify the mechanisms can consult stable in-repository sources: the
architecture map in \code{docs/ARCHITECTURE.md} and maintainer guide in
\code{AGENTS.md}; the Engineer round loop, session-roll, and stall guards in
\code{argus\_skill/engineer/}; the completion authority and progress classes in
\code{argus\_skill/reviewer/\_core.py} and its \code{reviewer\_schema.json}; the
Manager front door in \code{argus\_skill/manager/}; the 7$\times$24 scheduler in
\code{argus\_skill/life/supervisor/}; and the control-plane reliability record in
\code{docs/experiment/2026-07-12-completion-livelock/}, whose experiments verify
that four control-plane omissions no longer cause a stuck stage or a repeated
planning loop. The selected public results carry their own machine-readable
evidence bundles, \code{technical\_report/evidence/website\_results.json} and
\code{paper\_inventory.json}. This report cites repository paths and commit
lineage on \code{main} rather than any transient local filesystem location.
